\chapter{Non-Normal Linear Algebra}
\label{app:nonnormal}

This appendix collects the linear-algebra results that underpin the
non-Hermitian spectral theory used throughout the book. The central
theme is \emph{non-normality}: a matrix $A$ is normal if
$A A^\dagger = A^\dagger A$; Hermitian, unitary, and skew-Hermitian
matrices are all normal. A non-normal matrix can have a complete set
of eigenvectors, but they are not orthogonal, and the spectral behavior
under perturbation is qualitatively different from the normal case.
Non-Hermitian photonic operators---the Maxwell operator with complex
$\epsr$, the scattering matrix near an EP, the effective Hamiltonian of
a coupled-mode model---are generically non-normal.


%% ================================================================
\section{Left and right eigenvectors}
\label{sec:left-right}

For a matrix $A \in \mathbb{C}^{n \times n}$, the \emph{right
eigenvector} $|r_j\rangle$ and the \emph{left eigenvector}
$\langle l_j|$ associated with eigenvalue $\lambda_j$ satisfy
\begin{equation}
 A\,|r_j\rangle = \lambda_j\,|r_j\rangle, \qquad
 \langle l_j|\,A = \lambda_j\,\langle l_j|.
 \label{eq:left-right-def}
\end{equation}
Equivalently, $\langle l_j|$ is a right eigenvector of $A^\dagger$
with eigenvalue $\lambda_j^*$.

For a normal matrix, left and right eigenvectors coincide (up to
normalization):\ $|l_j\rangle = |r_j\rangle$. For a non-normal matrix,
they differ, and the angle between corresponding left and right
eigenvectors can be arbitrarily small---a situation that has dramatic
physical consequences (enhanced noise, sensitivity to perturbation,
mode non-orthogonality).


%% ================================================================
\section{Biorthogonality}
\label{sec:biorthogonality-app}

If $A$ has $n$ distinct eigenvalues, the left and right eigenvectors
form a \emph{biorthogonal} system:
\begin{equation}
 \langle l_i | r_j \rangle = \delta_{ij}\,N_j,
 \label{eq:biorthogonal}
\end{equation}
where $N_j = \langle l_j | r_j \rangle$ is a nonzero normalization
constant (which is generally complex for non-normal matrices). The
biorthogonal normalization convention sets $N_j = 1$ for all $j$ by
rescaling.

The spectral decomposition of $A$ takes the biorthogonal form
\begin{equation}
 A = \sum_{j=1}^n \lambda_j\,|r_j\rangle\langle l_j|,
 \label{eq:spectral-biorthogonal}
\end{equation}
and the resolvent is
\begin{equation}
 (z - A)^{-1} = \sum_{j=1}^n \frac{|r_j\rangle\langle l_j|}
 {z - \lambda_j}.
 \label{eq:resolvent-biorthogonal}
\end{equation}
These formulas generalize the Hermitian spectral theorem to the
non-normal case. They are the mathematical foundation for the
quasinormal-mode expansion (Chapter~\ref{ch:qnm}) and the
biorthogonal Berry phase (Chapter~\ref{ch:complex-bands}).


%% ================================================================
\section{Jordan chains and defective eigenvalues}
\label{sec:jordan-app}

When two or more eigenvalues coalesce, the eigenvectors may also
coalesce, producing a \emph{defective} eigenvalue---an exceptional
point (Chapter~\ref{ch:exceptional-points}).

\subsection{Jordan normal form}
\label{ssec:jordan-form}

At a defective eigenvalue $\lambda_0$ of algebraic multiplicity $m$ and
geometric multiplicity $g < m$, the matrix $A$ cannot be diagonalized.
Instead, the relevant block of the Jordan normal form is
\begin{equation}
 J_m(\lambda_0) = \begin{pmatrix}
 \lambda_0 & 1 & 0 & \cdots & 0 \\
 0 & \lambda_0 & 1 & \cdots & 0 \\
 \vdots & & \ddots & \ddots & \vdots \\
 0 & \cdots & 0 & \lambda_0 & 1 \\
 0 & \cdots & 0 & 0 & \lambda_0
 \end{pmatrix}
 \in \mathbb{C}^{m \times m}.
 \label{eq:jordan-block-app}
\end{equation}

\subsection{Generalized eigenvectors}
\label{ssec:generalized-evecs}

The Jordan chain associated with $J_m(\lambda_0)$ consists of vectors
$|v_0\rangle, |v_1\rangle, \ldots, |v_{m-1}\rangle$ satisfying
\begin{equation}
 (A - \lambda_0)\,|v_0\rangle = 0, \qquad
 (A - \lambda_0)\,|v_p\rangle = |v_{p-1}\rangle,
 \quad p = 1, \ldots, m-1.
 \label{eq:jordan-chain-app}
\end{equation}
Only $|v_0\rangle$ is a true eigenvector; the higher-order vectors
$|v_p\rangle$ are \emph{generalized eigenvectors}. At a second-order
EP ($m = 2$), there is one eigenvector and one generalized eigenvector;
the biorthogonal decomposition~\eqref{eq:spectral-biorthogonal} breaks
down and must be replaced by the Jordan decomposition.

\subsection{Resolvent near an EP}
\label{ssec:resolvent-ep}

Near a second-order EP at $\lambda_0$, the resolvent has a
second-order pole:
\begin{equation}
 (z - A)^{-1} \approx
 \frac{|v_1\rangle\langle w_0|}{(z - \lambda_0)^2}
 + \frac{|v_0\rangle\langle w_0| + |v_1\rangle\langle w_1|}
 {z - \lambda_0}
 + \text{regular},
 \label{eq:resolvent-ep}
\end{equation}
where $\langle w_0|$ and $\langle w_1|$ are the left Jordan-chain
vectors. The double pole produces algebraic (rather than exponential)
temporal growth---the hallmark of an EP in time-domain experiments and
the origin of the enhanced transient response.


%% ================================================================
\section{Petermann factor and excess noise}
\label{sec:petermann-app}

The Petermann factor quantifies the enhancement of spontaneous emission
(and, more generally, noise) in a non-normal system. For a mode with
right eigenvector $|r_j\rangle$ and left eigenvector $\langle l_j|$,
the Petermann factor is
\begin{equation}
 K_j = \frac{\langle r_j | r_j \rangle \;\langle l_j | l_j \rangle}
 {|\langle l_j | r_j \rangle|^2} \ge 1,
 \label{eq:petermann-def}
\end{equation}
with equality if and only if $|l_j\rangle = |r_j\rangle$ (the normal
case).

\subsection{Physical interpretation}
\label{ssec:petermann-physics-app}

In a laser, spontaneous emission couples to the lasing mode with an
amplitude proportional to $|\langle l_j | r_j \rangle|^{-1}$. When
the left and right eigenvectors are nearly orthogonal (as happens near
an EP), this coupling is enhanced, broadening the laser linewidth by
the factor $K_j$:
\begin{equation}
 \Delta\nu = K_j\,\Delta\nu_{\mathrm{ST}},
 \label{eq:petermann-linewidth-app}
\end{equation}
where $\Delta\nu_{\mathrm{ST}}$ is the Schawlow--Townes linewidth.

\subsection{Scaling near an EP}
\label{ssec:petermann-ep}

Near a second-order EP, the overlap $\langle l_j | r_j \rangle$
vanishes as $\sqrt{\epsilon}$ where $\epsilon$ is the distance from
the EP in parameter space. Therefore the Petermann factor diverges as
\begin{equation}
 K_j \sim \frac{1}{\epsilon}
 \label{eq:petermann-ep-scaling}
\end{equation}
near the EP. In common equal-resource sensing models, the associated
noise growth can offset the enhanced frequency splitting
$\Delta\omega \sim \sqrt{\epsilon}$ in a signal-to-noise proxy, as
discussed in Chapter~\ref{ch:devices}. The quantitative cancellation is
not universal; it depends on the measurement protocol, noise channels,
and normalization of the sensor response.

For an $n$-th order EP, analogous power-law divergences appear, but the
resulting precision bound must be evaluated from the full noise model
rather than inferred from the eigenvalue splitting alone.


%% ================================================================
\section{Phase rigidity}
\label{sec:phase-rigidity}

The \emph{phase rigidity} of a mode is the inverse of the Petermann
factor:
\begin{equation}
 r_j = \frac{|\langle l_j | r_j \rangle|^2}
 {\langle r_j | r_j \rangle \;\langle l_j | l_j \rangle}
 = K_j^{-1} \in [0, 1].
 \label{eq:phase-rigidity-def}
\end{equation}
A phase rigidity of $r_j = 1$ indicates a normal (Hermitian-like) mode;
$r_j \to 0$ indicates an EP where the mode becomes maximally
non-orthogonal.


%% ================================================================
\section{Pseudospectra}
\label{sec:pseudospectra}

The spectrum of a non-normal matrix can be misleading: eigenvalues may
be highly sensitive to perturbations even when they are well separated.
The \emph{$\epsilon$-pseudospectrum} quantifies this sensitivity.

\subsection{Definition}
\label{ssec:pseudospec-def}

The $\epsilon$-pseudospectrum of $A$ is the set of complex numbers
$z$ such that $z$ is an eigenvalue of $A + E$ for some perturbation
$\|E\| \le \epsilon$:
\begin{equation}
 \sigma_\epsilon(A) = \{z \in \mathbb{C} :
 \|(z - A)^{-1}\| \ge \epsilon^{-1}\}.
 \label{eq:pseudospectrum}
\end{equation}
For a normal matrix, the $\epsilon$-pseudospectrum is the union of
$\epsilon$-disks around each eigenvalue. For a non-normal matrix, the
pseudospectrum can extend far beyond these disks, indicating that
eigenvalues can move large distances under small perturbations.

\subsection{Physical relevance}
\label{ssec:pseudospec-physics}

Pseudospectra are relevant whenever the system is subject to
perturbations---fabrication disorder, environmental fluctuations,
or numerical discretization errors. For a photonic system operating
near an EP, the pseudospectrum extends over a large region of the
complex frequency plane, meaning that the eigenfrequencies are
effectively unpredictable to within a band set by the perturbation
strength. This is the spectral manifestation of the Petermann-factor
divergence and the reason why EP-based devices require careful
noise analysis (Chapter~\ref{ch:devices}).


%% ================================================================
\section{Matrix functions and the non-Hermitian exponential}
\label{sec:matrix-exp}

The time evolution of a system governed by $i\,\partial_t |\psi\rangle
= A\,|\psi\rangle$ is
\begin{equation}
 |\psi(t)\rangle = e^{-iAt}\,|\psi(0)\rangle.
 \label{eq:time-evolution}
\end{equation}
For a normal $A$, the evolution is a superposition of oscillating and
decaying exponentials, and $\|e^{-iAt}\|$ is controlled by the largest
imaginary part of the eigenvalues. For a non-normal $A$, the situation
is richer: the Jordan chain at a defective eigenvalue contributes
\emph{polynomial} growth factors
\begin{equation}
 e^{-iJ_m(\lambda_0)t} = e^{-i\lambda_0 t}
 \begin{pmatrix}
 1 & -it & (-it)^2/2! & \cdots \\
 0 & 1 & -it & \cdots \\
 \vdots & & \ddots & \ddots \\
 0 & \cdots & 0 & 1
 \end{pmatrix},
 \label{eq:jordan-exponential}
\end{equation}
so the amplitude can grow as $t^{m-1}$ even when the eigenvalue has
a negative imaginary part (i.e., the mode is formally decaying). This
\emph{transient growth} is a hallmark of non-normality and is
observable as a power-law ring-up before the eventual exponential
decay in photonic resonators near an EP.

Additionally, even in the absence of defective eigenvalues, a highly
non-normal matrix can produce large transient growth---the so-called
\emph{pseudospectral} transient, where the norm $\|e^{-iAt}\|$ overshoots
its long-time asymptote by a factor that can be estimated from the
pseudospectrum. This effect is important for FDTD simulations of
non-Hermitian photonic systems (Section~\ref{sec:fdtd}), where a
transient overshoot can trigger numerical instability before the
physical decay sets in.


%% ================================================================
\section{Examples}
\label{sec:nonnormal-examples}

\subsection{Jordan decomposition of a $2\times 2$ defective matrix}
\label{ssec:jordan-2x2}

The simplest defective matrix arises at an exceptional point of a
$2\times 2$ non-Hermitian Hamiltonian. Consider
\begin{equation}
 A = \begin{pmatrix} \lambda_0 & 1 \\ 0 & \lambda_0 \end{pmatrix},
 \label{eq:jordan-2x2-example}
\end{equation}
which is already in Jordan normal form $J_2(\lambda_0)$. This matrix
has a single eigenvalue $\lambda_0$ with algebraic multiplicity~2 but
geometric multiplicity~1: there is only one linearly independent
eigenvector, $\mathbf{v}_1 = (1, 0)^\transpose$. The second vector in
the Jordan chain is the \emph{generalized eigenvector} $\mathbf{v}_2$
satisfying $(A - \lambda_0 I)\mathbf{v}_2 = \mathbf{v}_1$, which gives
$\mathbf{v}_2 = (0, 1)^\transpose$.

The time evolution is
\begin{equation}
 e^{-iAt} = e^{-i\lambda_0 t}
 \begin{pmatrix} 1 & -it \\ 0 & 1 \end{pmatrix}.
 \label{eq:jordan-2x2-evolution}
\end{equation}
Even if $\Im\lambda_0 < 0$ (the mode decays), the off-diagonal element
$-it$ grows linearly before the exponential decay dominates. The
maximum of $\|e^{-iAt}\|$ over time occurs at
$t^* \sim 1/|\Im\lambda_0|$, with an envelope of order
$1/(e|\Im\lambda_0|)$ before constant prefactors from the chosen norm
are included. The transient amplification therefore diverges as the
eigenvalue approaches the real axis, but the numerical prefactor is
norm dependent.

\subsection{Schur decomposition and its use}
\label{ssec:schur-decomp}

Every square matrix $A$ admits a Schur decomposition $A = Q T Q^{-1}$,
where $Q$ is unitary and $T$ is upper triangular with the eigenvalues on
the diagonal. Unlike the Jordan form, which requires exact knowledge of
the eigenvector structure, the Schur form is numerically stable and can
be computed by standard algorithms (QR iteration).

The Schur form reveals the \emph{departure from normality}
$\Delta_F = \|A\|_F^2 - \sum_i |\lambda_i|^2$, where $\|\cdot\|_F$
is the Frobenius norm. For a normal matrix, $\Delta_F = 0$ (the Schur
form is diagonal); for a non-normal matrix, $\Delta_F > 0$ measures the
energy stored in the off-diagonal elements of $T$, which is the source
of transient amplification and pseudospectral broadening.

For photonic applications, the Schur decomposition is preferable to the
Jordan form for three reasons. First, it is numerically robust: no
special treatment is needed for nearly defective matrices. Second, the
unitary factor $Q$ preserves norms, making it easy to compute
$\|e^{-iAt}\|$ from $\|e^{-iTt}\|$. Third, the off-diagonal elements
of $T$ give a direct and intuitive measure of the non-Hermitian coupling
between modes---elements that are absent in normal systems.

\subsection{Kreiss constant and transient bounds}
\label{ssec:kreiss}

The maximum transient growth $\sup_{t \ge 0}\|e^{-iAt}\|$ can be
bounded using the Kreiss matrix theorem. The \emph{Kreiss constant}
\begin{equation}
 \mathcal{K}(A) = \sup_{\Im z > \alpha_{\max}}
 (\Im z - \alpha_{\max})\,\|(z - A)^{-1}\|,
 \label{eq:kreiss-constant}
\end{equation}
where $\alpha_{\max} = \max_i \Im\lambda_i$ is the maximum eigenvalue
imaginary part, satisfies the bounds
\begin{equation}
 \mathcal{K}(A) \le \sup_{t \ge 0} \|e^{-iAt}\|\,e^{\alpha_{\max}t}
 \le e\,n\,\mathcal{K}(A),
 \label{eq:kreiss-bounds}
\end{equation}
where $n$ is the matrix dimension. The Kreiss constant is computable
from the resolvent norm (essentially the pseudospectrum) and provides
a practical upper bound on the transient growth without requiring
explicit time integration.

For the $2\times 2$ PT-symmetric Hamiltonian near the EP
($\gamma/\kappa = 0.99$), the Kreiss constant is
$\mathcal{K} \approx 7.1$, consistent with the transient amplification
factor of $7.09$ computed directly from~\eqref{eq:transient-max}. For
larger photonic systems, the Kreiss constant can be estimated from the
pseudospectrum using standard numerical linear algebra packages.

\subsection{Non-normal perturbation theory}
\label{ssec:nonnormal-perturbation}

For a non-normal matrix $A$ with a defective eigenvalue $\lambda_0$ of
algebraic multiplicity $m$, a generic perturbation $\epsilon B$ splits
the eigenvalue into $m$ branches:
\begin{equation}
 \lambda_j(\epsilon) = \lambda_0
 + \epsilon^{1/m}\,c_j + O(\epsilon^{2/m}),
 \qquad j = 1, \ldots, m,
 \label{eq:puiseux-series}
\end{equation}
where $c_j$ are the $m$-th roots of a coupling coefficient that depends
on the perturbation matrix $B$ and the Jordan chain vectors. This
Puiseux series expansion is the mathematical basis for the enhanced
spectral response at exceptional points: a perturbation of order
$\epsilon$ produces an eigenvalue splitting of order
$\epsilon^{1/m}$, which for small $\epsilon$ is much larger than the
linear splitting $\sim\epsilon$ of a non-defective eigenvalue.

However, the Puiseux expansion is \emph{not} isotropic in the
perturbation: the coefficient $c_j$ depends on the \emph{direction} of
the perturbation in parameter space, not just its magnitude. For
certain perturbation directions (those orthogonal to the Jordan chain
coupling), the splitting can be much smaller than $\epsilon^{1/m}$ or
even vanish. This anisotropy is important for EP-based sensing: the
enhanced response occurs only for perturbations that couple to the
Jordan chain, not for arbitrary environmental noise.
